
\subsection{%
Corollaries of the Main Theorem
}

First we derive the equivariant cohomology version of Theorem~\ref{thm:main1}.
In addition to translating Nakada's colored hook formula~\cite[Corollary~7.2]{N1} 
from the context of roots to the context of $d$-complete posets, 
the following corollary gives a skew generalization of it.

\begin{coro}
\label{cor:main2}
Let $P$ be a $d$-complete poset with $d$-complete coloring $c : P \to I$ 
and $F$ an order filter of $P$.
Let $\vecta = (a_i)_{i \in I}$ be indeterminates.
We put $a(p) = a_{c(p)}$ \textup{(}$p \in P$\textup{)} 
and define a linear polynomial $\vecta \langle H_P(u) \rangle$ 
as follows:
\begin{enumerate}[(i)]
\item\label{coro5.4_i} %[(i)]
If $u$ is not the top of any $d_k$-interval, then we define
$$
\vecta \langle H_P(u) \rangle = \sum_{w \le u} a_{c(w)}.
$$
\item\label{coro5.4_ii} %[(ii)]
If $u$ is the top of a $d_k$-interval $[v,u]$, then we define
$$
\vecta \langle H_P(u) \rangle
 =
\vecta \langle H_P(x) \rangle + \vecta \langle H_P(y) \rangle
- \vecta \langle H_P(v) \rangle,
$$
where $x$ and $y$ are the sides of $[v,u]$.
\end{enumerate}
Then we have
\begin{multline}
\label{eq:main2}
\sum_{(q_1, \dots, q_n)}
 \frac{ 1 }
 { a(q_1) (a(q_1)+a(q_2)) \cdots (a(q_1)+\dots+a(q_n)) }
\\
 =
\sum_{D \in \EE_P(F)}
 \prod_{v \in P \setminus D}
 \frac{ 1 }
 { \vecta \langle H_P(v) \rangle },
\end{multline}
where $n = \# (P\setminus F)$ and 
the summation is taken over all linear extensions of $P \setminus F$, 
i.e. all labelings of the elements of $P \setminus F$ with $q_1, \dots, q_n$ so that 
 $q_i < q_j$ in $P \setminus F$ implies $i < j$.
\end{coro}

\begin{proof}
Any $(P \setminus F)$-partition $\sigma \in \AP(P \setminus F)$ is determined by 
a nonnegative integer $r\leq n$, an increasing sequence $i_1 < \cdots < i_r$ of positive integers 
and an increasing sequence $F \subset F_0 \subsetneq F_1 \subsetneq \cdots \subsetneq F_r = P$ 
of order filters of $P$, by the condition
$$
\sigma(x)
 =
\begin{cases}
 0 &\text{if $x \in F_0 \setminus F$,} \\
 i_k &\text{if $x \in F_k \setminus F_{k-1}$ and $1 \le k \le r$.}
\end{cases}
$$
Hence we have
\begin{multline*}
\sum_{\sigma \in \AP(P \setminus F)} \vectz^\sigma
\\
\begin{aligned}
&
=
\sum_{F \subset F_0 \subsetneq F_1 \subsetneq \cdots \subsetneq F_r = P}
\sum_{0 < i_1 < \cdots < i_r}
 \vectz[F_1 \setminus F_0]^{i_1} \vectz[F_2 \setminus F_1]^{i_2} \cdots \vectz[P \setminus F_{r-1}]^{i_r}
\\
&
=
\sum_{F \subset F_0 \subsetneq F_1 \subsetneq \cdots \subsetneq F_r = P}
 \frac{ \vectz[P \setminus F_0] }
 { 1 - \vectz[P \setminus F_0] }
 \frac{ \vectz[P \setminus F_1] }
 { 1 - \vectz[P \setminus F_1] }
 \cdots
 \frac{ \vectz[P \setminus F_{r-1}] }
 { 1 - \vectz[P \setminus F_{r-1}] }.
\end{aligned}
\end{multline*}
Now by using Theorem~\ref{thm:main1}, we have
\begin{multline*}
\sum_{F \subset F_0 \subsetneq F_1 \subsetneq \cdots \subsetneq F_r = P}
 \frac{ \vectz[P \setminus F_0] }
 { 1 - \vectz[P \setminus F_0] }
 \frac{ \vectz[P \setminus F_1] }
 { 1 - \vectz[P \setminus F_1] }
 \cdots
 \frac{ \vectz[P \setminus F_{r-1}] }
 { 1 - \vectz[P \setminus F_{r-1}] }
\\
 =
\sum_{D \in \EE_P(F)}
 \frac{ \prod_{v \in B(D)} \vectz[H_P(v)] }
 { \prod_{v \in P \setminus D} ( 1 - \vectz[H_P(v)] ) }.
\end{multline*}
By substituting $z_i = t^{a_i}$ ($i \in I$) and multiplying the both sides by $(1-t)^n$, 
and then by taking the limit $t \to 1$, we obtain
\begin{multline*}
\sum_{F = F_0 \subsetneq F_1 \subsetneq \cdots \subsetneq F_n = P}
 \frac{ 1 }
 { \vecta \langle P \setminus F_0 \rangle \vecta \langle P \setminus F_1 \rangle 
 \cdots \vecta \langle P \setminus F_{n-1} \rangle
 }
\\
 =
\sum_{D \in \EE_P(F)}
 \prod_{v \in P \setminus D}
 \frac{ 1 }
 { \vecta \langle H_P(v) \rangle },
\end{multline*}
where the summation on the left hand side is taken over all increasing sequences 
$F = F_0 \subsetneq F_1 \subsetneq \cdots \subsetneq F_n = P$ of order filters of length $n$, 
and $\vecta \langle D \rangle = \sum_{p \in D} a_{c(p)}$ for a subset $D \subset P$.
Such increasing sequences of order filters are in one-to-one correspondence 
with linear extensions $(q_1, \ldots, q_n)$ of $P \setminus F$ by the relation
$$
F_k = F \cup \{ q_n, \ldots, q_{n-k+1} \}
\quad(0 \le k \le n).
$$
Hence we obtain the desired result.
\end{proof}

\begin{rema}
Corollary~\ref{cor:main2} can be proved by using the Billey formula~\cite[Theorem~4]{B} 
and the Chevalley formula~\cite[Theorem~11.1.7 and Corollary~11.3.17]{Ku} 
for the equivariant cohomology along the same line as Theorem~\ref{thm:main1}.
\end{rema}

By specializing $z_i = q$ for all $i \in I$ in~\eqref{eq:main1}, 
and $a_i = 1$ for all $i \in I$ in~\eqref{eq:main2}, we~obtain

\begin{coro}
\label{cor:main3}
Let $P$ be a $d$-complete poset and $F$ an order filter of $P$.
We define the hook length $h_P(u)$ at $u \in P$ as follows:
\begin{enumerate}[(i)]
\item\label{coro5.6_i} %[(i)]
If $u$ is not the top of any $d_k$-interval, then we define $h_P(u) = \# \{ w \in P : w \le u \}$.
\item\label{coro5.6_ii} %[(ii)]
If $u$ is the top of a $d_k$-interval $[v,u]$, then we define 
$h_P(u) = h_P(x) + h_P(y) - h_P(v)$, 
where $x$ and $y$ are the sides of $[v,u]$.
\end{enumerate}
Then we have
\begin{enumerate}[(a)]
\item\label{coro5.6_a} %[(a)]
The univariate generating function of $(P \setminus F)$-partitions is given by
\begin{equation}
\label{eq:main31}
\sum_{\sigma \in \AP(P \setminus F)}
 q^{|\sigma|}
 =
\sum_{D \in \EE_P(F)}
\frac{ \prod_{v \in B(D)} q^{h_P(v)} }
         { \prod_{v \in P \setminus D} (1-q^{h_P(v)}) }.
\end{equation}
\item\label{coro5.6_b} %[(b)]
The number of linear extensions of $P \setminus F$ is given by
\begin{equation}
\label{eq:main32}
n! \sum_{D \in \EE_P(F)} \prod_{v \in P \setminus D} \frac{1}{h_P(v)},
\end{equation}
where $n = \# (P \setminus F)$.
\end{enumerate}
\end{coro}

If $P = D(\lambda)$ and $F = D(\mu)$ are shapes corresponding to partitions $\lambda \supset \mu$,
Equations~\eqref{eq:main31} and \eqref{eq:main32} reduce to 
Morales--Pak--Panova's $q$-hook formula~\cite[Corollary~6.17]{MPP} and Naruse's hook formula~\cite{Naruse} 
respectively.
The trace generating function of reverse plane partitions of skew shape~\cite[Corollary~6.20]{MPP}
is obtained from Theorem~\ref{thm:main1} by specializing
$$
z_i = \begin{cases}
 tq &\text{if $i$ is the color of the maximum element of $D(\lambda)$,} \\
 q &\text{otherwise.}
\end{cases}
$$

\begin{rema}
\label{rem:heap3}
Theorem~\ref{thm:main1} and its corollaries hold for heaps $H(w)$ associated 
to dominant minuscule elements $w$ in any symmetrizable Kac--Moody Weyl groups, 
after suitable modifications are made.
See Remarks~\ref{rem:heap1} and~\ref{rem:heap2}.
\end{rema}

\subsection{%
Example
}

\begin{exam}
\label{ex:dc}
Let $P = S(3,2,1) \supset F = S(2)$ be the shifted shapes corresponding 
to strict partitions $(3,2,1)$ and $(2)$.
If we regard $P$ as a $d$-complete poset with a $d$-complete coloring $c : P \to \{ 0, 0', 1, 2 \}$ 
given in Example~\ref{ex:shifted_hc}, 
then the hook monomials in 
${{\vectz}} = (z_0, z_{0'}, z_1, z_2)$ are given by
\begin{alignat*}{3}
{{\vectz}}[ H_P(1,1) ] &= z_0 z_{0'} z_1^2 z_2,
&\quad
{{\vectz}}[ H_P(1,2) ] &= z_0 z_{0'} z_1 z_2,
&\quad
{{\vectz}}[ H_P(1,3) ] &= z_0 z_1 z_2,
\\
&
&
{{\vectz}}[ H_P(2,2) ] &= z_0 z_{0'} z_1,
&\quad
{{\vectz}}[ H_P(2,3) ] &= z_0 z_1,
\\
&
&
&
&{{\vectz}}[ H_P(3,3) ] &= z_0.
\end{alignat*}
Since we have
$$
\setlength{\unitlength}{1.1pt}
\mathcal{E}_P(F)
 =
\left\{
\raisebox{-14pt}{
\begin{picture}(30,30)
\fboxsep=0mm
\put(0,20){\colorbox[gray]{0.7}{\makebox(10,10){}}}
\put(10,20){\colorbox[gray]{0.7}{\makebox(10,10){}}}
\put(0,30){\line(1,0){30}}
\put(0,20){\line(1,0){30}}
\put(10,10){\line(1,0){20}}
\put(20,0){\line(1,0){10}}
\put(0,20){\line(0,1){10}}
\put(10,10){\line(0,1){20}}
\put(20,0){\line(0,1){30}}
\put(30,0){\line(0,1){30}}
\end{picture}
},
\ 
\raisebox{-14pt}{
\begin{picture}(30,30)
\fboxsep=0mm
\put(0,20){\colorbox[gray]{0.7}{\makebox(10,10){}}}
\put(20,10){\colorbox[gray]{0.7}{\makebox(10,10){}}}
\put(0,30){\line(1,0){30}}
\put(0,20){\line(1,0){30}}
\put(10,10){\line(1,0){20}}
\put(20,0){\line(1,0){10}}
\put(0,20){\line(0,1){10}}
\put(10,10){\line(0,1){20}}
\put(20,0){\line(0,1){30}}
\put(30,0){\line(0,1){30}}
\put(10,20){\makebox(10,10){$\times$}}
\end{picture}
}
\right\},
$$
we apply Theorem~\ref{thm:main1} to obtain
\begin{align}
\label{eq:ex1}
\sum_{\pi \in \AP(P \setminus F)} {\vectz}^\pi
 &=
\frac{ 1 }
 { (1 - z_0 z_1 z_2) (1 - z_0 z_{0'} z_1) (1 - z_0 z_1) (1 - z_0) }
\\
 &\quad
+
\frac{ z_0 z_{0'} z_1 z_2 }
 { (1 - z_0 z_{0'} z_1 z_2) (1 - z_0 z_1 z_2) (1 - z_0 z_{0'} z_1) (1 - z_0) }
\notag
\\
 &=
\frac{ 1 - z_0^2 z_{0'} z_1^2 z_2 }
 { (1 - z_0 z_{0'} z_1 z_2) (1 - z_0 z_1 z_2) (1 - z_0 z_{0'} z_1) (1 - z_0 z_1) (1 - z_0) },
\notag
\end{align}
where
$$
{\vectz}^\pi
 =
z_0^{\pi(1,1) + \pi(3,3)} z_{0'}^{\pi(2,2)} z_1^{\pi(1,2) + \pi(2,3)} z_2^{\pi(1,3)}.
$$
\end{exam}

\begin{exam}
\label{ex:heap}
Let $P = S(3,2,1) \supset F = S(2)$ be the same as in Example~\ref{ex:dc}.
If we regard $P$ as the heap $H(w_{(3,2,1)})$ for the Weyl group of type $B_3$ 
(see Example~\ref{ex:shifted-B}),
then the hook monomials in $\overline{\vectz} = (z_0, z_1, z_2)$ are given by
\begin{alignat*}{3}
\overline{\vectz}[ H'_P(1,1) ] &= z_0 z_1 z_2,
&\quad
\overline{\vectz}[ H'_P(1,2) ] &= z_0^2 z_1^2 z_2,
&\quad
\overline{\vectz}[ H'_P(1,3) ] &= z_0^2 z_1 z_2,
\\
&
&
\overline{\vectz}[ H'_P(2,2) ] &= z_0 z_1,
&\quad
\overline{\vectz}[ H'_P(2,3) ] &= z_0^2 z_1,
\\
&
&
&
&
\overline{\vectz}[ H'_P(3,3) ] &= z_0.
\end{alignat*}
Since we have
$$
\setlength{\unitlength}{1.1pt}
\mathcal{E}_P(F)
 =
\left\{
\raisebox{-14pt}{
\begin{picture}(30,30)
\fboxsep=0mm
\put(0,20){\colorbox[gray]{0.7}{\makebox(10,10){}}}
\put(10,20){\colorbox[gray]{0.7}{\makebox(10,10){}}}
\put(0,30){\line(1,0){30}}
\put(0,20){\line(1,0){30}}
\put(10,10){\line(1,0){20}}
\put(20,0){\line(1,0){10}}
\put(0,20){\line(0,1){10}}
\put(10,10){\line(0,1){20}}
\put(20,0){\line(0,1){30}}
\put(30,0){\line(0,1){30}}
\end{picture}
},
\ 
\raisebox{-14pt}{
\begin{picture}(30,30)
\fboxsep=0mm
\put(0,20){\colorbox[gray]{0.7}{\makebox(10,10){}}}
\put(20,10){\colorbox[gray]{0.7}{\makebox(10,10){}}}
\put(0,30){\line(1,0){30}}
\put(0,20){\line(1,0){30}}
\put(10,10){\line(1,0){20}}
\put(20,0){\line(1,0){10}}
\put(0,20){\line(0,1){10}}
\put(10,10){\line(0,1){20}}
\put(20,0){\line(0,1){30}}
\put(30,0){\line(0,1){30}}
\put(10,20){\makebox(10,10){$\times$}}
\end{picture}
},
\ 
\raisebox{-14pt}{
\begin{picture}(30,30)
\fboxsep=0mm
\put(10,10){\colorbox[gray]{0.7}{\makebox(10,10){}}}
\put(20,10){\colorbox[gray]{0.7}{\makebox(10,10){}}}
\put(0,30){\line(1,0){30}}
\put(0,20){\line(1,0){30}}
\put(10,10){\line(1,0){20}}
\put(20,0){\line(1,0){10}}
\put(0,20){\line(0,1){10}}
\put(10,10){\line(0,1){20}}
\put(20,0){\line(0,1){30}}
\put(30,0){\line(0,1){30}}
\put(0,20){\makebox(10,10){$\times$}}
\end{picture}
}
\right\},
$$
we apply a heap version of Theorem~\ref{thm:main1} to obtain
\begin{align}
\label{eq:ex2}
\sum_{\pi \in \AP(P \setminus F)} \overline{\vectz}^\pi
 &=
\frac{ 1 }
 { (1 - z_0^2 z_1 z_2) (1 - z_0 z_1) (1 - z_0^2 z_1) (1 - z_0) }
\\
 &\quad
+
\frac{ z_0^2 z_1^2 z_2 }
 { (1 - z_0^2 z_1^2 z_2) (1 - z_0^2 z_1 z_2) (1 - z_0 z_1) (1 - z_0) }
\notag
\\
 &\quad
+
\frac{ z_0 z_1 z_2 }
 { (1 - z_0 z_1 z_2) (1 - z_0^2 z_1^2 z_2) (1 - z_0^2 z_1 z_2) (1 - z_0) }
\notag
\\
 &=
\frac{ 1 - z_0^3 z_1^2 z_2 }
 { (1 - z_0 z_1 z_2) (1 - z_0^2 z_1 z_2) (1 - z_0 z_1) (1 - z_0^2 z_1) (1 - z_0) },
\notag
\end{align}
where
$$
\overline{\vectz}^\pi
 =
z_0^{\pi(1,1) + \pi(2,2) + \pi(3,3)} z_1^{\pi(1,2) + \pi(2,3)} z_2^{\pi(1,3)}.
$$
Note that Equation~\eqref{eq:ex2} is obtained from~\eqref{eq:ex1} by putting $z_{0'} = z_0$.
\end{exam}

